\documentclass[../Thesis-IJspeert.tex]{subfiles}

\begin{document}

\graphicspath{ {"Cavity-Compatible Trapping Architectures/figs/"} }
\pgfplotsset{table/search path={"Cavity-Compatible Trapping Architectures/data/"}}

\chapter{Cavity-Compatible Trapping Architectures}
\addtocontents{toc}{\vskip-6pt\par\noindent\protect\textcolor{gray75}{\protect\rule{\textwidth}{0.5pt}}\par}
\label{chap:CavityCompatibleTrappingArchitectures}

\section{CAVITY MOUNTING FINGER DESIGN}
\subsection{Overview}
\subsection{Structural analysis}
In the current setup for cavity-assisted photon generation \cite{Barrett2019}, atoms are loaded stochastically into the cavity mode via an atomic fountain. However, the cavity used in that experiment (with a geometry as shown in \autoref{fig:cavity}) is not compatible with holographic optical tweezers, since its mirrors are relatively large ($\diameter \SI{1.5}{\milli\meter}$) in comparison to the length of the cavity (\SI{339}{\micro\meter}). These dimensions imply that the full accessible angle of this cavity is \ang{25.5}, whereas the trapping light would propagate through the glass cell at nearly the full acceptance angle of the high numerical aperture lens, which is \ang{73.7}. A new cavity is under construction which consists of two truncated square pyramids of silica. An important difference is that the surfaces of truncation are much smaller ($\SI{300}{\micro\meter}\times\SI{300}{\micro\meter}$). This implies that for a similar cavity length ($\sim\SI{300}{\micro\meter}$), the new design would feature a full accessible angle of \ang{90}, which is compatible with the current optical tweezers. Another difference is the method of cavity mirror fabrication. In the current design, the macroscopic silica substrates are super-polished, which means that these mirrors can have exceptionally low losses ($\sim\SI{2}{ppm}$). The disadvantage of this technique is that it results in a relatively large radius of curvature (\SI{50}{\milli\meter}) and mirror diameter. An alternative approach involves the use of intense \ce{CO2} laser pulses to ablate Gaussian micro-dimples onto the end facets of optical fibres \cite{Hunger2010}, which yields a much smaller radius of curvature (\SI{200}{\micro\meter} within a range varying from \SI{50}{\micro\meter} to \SI{1}{\milli\meter}) that can enable a smaller mode volume and therefore higher cooperativity. However, fibre-based optical microcavities can be difficult to stabilise due to the limited mass of the fibre and the fact that the cavity facets are always under some tension from the rest of the structure \cite{Janitz2017}. Another complication is the limited spatial mode matching between the cavity mode and the input/output modes of the fibre. In the new cavity design, the ablation is performed on the truncated surfaces of the macroscopic silica bases for increased stability and improved mode matching, whilst incorporating the advantages of an ablated mirror surface.

\begin{figure}[t]
	\centering
	\begin{tikzpicture}[remember picture]
		\node[anchor=south west,inner sep=0] (image1) at (0,0) {\includegraphics[width=1\textwidth]{{"finger"}.pdf}};
		
		% left four
		\begin{scope}[x={(image1.south east)},y={(image1.north west)}]
			% draw a grid
			%\draw[help lines,xstep=.1,ystep=.1,overlay] (0,0) grid (1,1);
			% draw ticks
			%\foreach \x in {0,1,...,9} { \node [anchor=north,overlay] at (\x/10,0) {0.\x}; }
			%\foreach \y in {0,1,...,9} { \node [anchor=east,overlay] at (0,\y/10) {0.\y}; }
			
			% labels
			\node[markgrijs, align=center, rotate=0] at (0.05,0.9) {\footnotesize\sc (a)};
			\node[markgrijs, align=center, rotate=0] at (0.05,0.42) {\footnotesize\sc (b)};
			\node[markgrijs, align=center, rotate=0] at (0.67,0.42) {\footnotesize\sc (c)};
			
		\end{scope}
		
	\end{tikzpicture}
	\caption[Cavity mounting finger and integration into the vacuum system]{Cavity mounting finger and integration into the vacuum system. (\textsc{a}) The cavity mounting finger, which is designed to secure the cavity module with six screws. A small mount inserted from the side holds four mirrors in a pyramidal configuration to form a pyramid MOT from which to load the dipole trap. The insets on the right illustrate the assembly of the pyramid MOT mirrors within the mount. The mounting finger features multiple openings to provide optical access for (horizontal) cooling and (vertical) trapping beams. (\textsc{b}) Integration of the cavity mounting finger into the vacuum system. A cuboid glass cell encloses the mounting finger, which is bolted to a flange attached to the vacuum chamber. This flange also secures the mounting finger with six smaller bolts (see inset). (\textsc{c}) A partial view of the mounting finger inside the glass cell, showing a pair of high-NA lenses that focus and recollimate the trapping beam (shown in red).}
	\label{fingeroverview} 
\end{figure}


\begin{figure}[!t]
	%\vspace{0.7em}
	\centering
	\begin{tikzpicture}
	
	\begin{groupplot}[group style={group size=1 by 3, vertical sep=1cm},xmin=0,ymin=0,height=5.5cm,width=13.5cm,no markers]
	
	\iffalse
	\nextgroupplot
	[
	ylabel={},
	xmin=0, xmax=1,
	ymin=0, ymax=1,
	xticklabels={,\textcolor{white}{0},},
	yticklabels={,,},
	%legend pos=north west,
	ymajorgrids=true,
	xmajorgrids=true,
	grid style=dashed,
	%axis lines=none,        % No axis lines
	%tick style=none,        % No ticks
	%xtick=\empty,           % No x-axis ticks
	%ytick=\empty,           % No y-axis ticks
	xlabel={},              % No x-axis label
	ylabel={},              % No y-axis label
	%enlargelimits=true,     % Keeps the plot within the bounds
	]
	
	\addplot graphics [
	xmin=0, xmax=1,  % Stretch to fit the x-axis limits
	ymin=0, ymax=1    % Stretch to fit the y-axis limits
	] {FINALFINGER.png};
	\fi
	
	
	\nextgroupplot
	[
	ymode=log,
	ymin=1e-0,
	ymax=1e4,
	%xlabel={Frequency},
	%ylabel={Max. linear acceleration},
	xmin=0, xmax=10,
	xticklabels={,,},
	%ymin=-0.1, ymax=1.1,
	%ymin=0, ymax=120,
	%xtick={0,20,40,60,80,100},
	%ytick={0,20,40,60,80,100,120},
	%legend pos=north west,
	ymajorgrids=true,
	xmajorgrids=true,
	grid style=dashed,
	legend cell align={left},
	legend style={
		at={(1, 1)},  % Adjust (x, y) to move it
		anchor=north east,  % Attach the legend to its top-right corner
		draw=none,           % Optional: Remove border around the legend
		fill=none,
		align=left
	}
	]
	
	%\node[align=center] at (axis cs: 0.267, 3500) {\footnotesize A};
	\node[align=center] at (axis cs: 1.105, 3500) {\footnotesize\sc b};
	\node[align=center] at (axis cs: 1.254+0.2, 11) {\footnotesize\sc c};
	\node[align=center] at (axis cs: 2.940, 2000) {\footnotesize\sc d};
	%\node[align=center] at (axis cs: 3.024, 2500) {\footnotesize E};
	\node[align=center] at (axis cs: 3.893, 1500) {\footnotesize\sc f};
	%\node[align=center] at (axis cs: 6.643, 300) {\footnotesize G};
	\node[align=center] at (axis cs: 6.935, 1600) {\footnotesize\sc h};
	
	
	\addplot[
	color=markrood,
	name path=xpointone,
	] table [col sep = comma, x expr = \thisrowno{0}/1000, y index = 1]{M_x_01.csv};
	
	\addplot[
	color=rainbow2of8,
	name path=xone,
	] table [col sep = comma, x expr = \thisrowno{0}/1000, y index = 1]{M_x_1.csv};
	
	\addplot[
	color=rainbow4of8,
	name path=xfive,
	] table [col sep = comma, x expr = \thisrowno{0}/1000, y index = 1]{M_x_5.csv};
	
	\addplot+[draw=none,domain=0:10,name path=xzero] {0.01}; 
	
	\addplot[color=markrood, opacity=0.5] fill between[
	of = xone and xpointone,
	]; 
	
	\addplot[color=rainbow2of8, opacity=0.5] fill between[
	of = xfive and xone,
	]; 
	
	\addplot[color=rainbow4of8, opacity=0.5] fill between[
	of = xzero and xfive,
	]; 
	
	\node at (rel axis cs: 0.52, 0.65) [anchor=center] {\includegraphics[width=3.5cm]{MARWANFINGERtrans.png}};
	
	\addlegendentry{\footnotesize \hspace{0.3em} $G=\SI{0.1}{\percent}$}
	\addlegendentry{\footnotesize \hspace{0.3em} $G=\SI{1}{\percent}$}
	\addlegendentry{\footnotesize \hspace{0.3em} $G=\SI{5}{\percent}$}
	
	\draw[thick, <->] (rel axis cs: 0.64, 0.8) --++ (150:0.5cm);
	\node[align=left] at (rel axis cs: 0.632, 0.88) {\footnotesize $x$};
	
	\nextgroupplot
	[
	ymode=log,
	ymin=1e-0,
	ymax=1e3,
	%ymin=1e-2, ymax=1e0,
	%xlabel={Frequency},
	ylabel={Maximum linear acceleration magnitude (\si{\milli\meter\per\second\squared})},
	xmin=0, xmax=10,
	xticklabels={,,},
	%ymin=-0.1, ymax=1.1,
	%ymin=0, ymax=120,
	%xtick={0,20,40,60,80,100},
	%ytick={0,20,40,60,80,100,120},
	%legend pos=north west,
	ymajorgrids=true,
	xmajorgrids=true,
	grid style=dashed,
	legend cell align={left},
	legend style={
		at={(1, 1)},  % Adjust (x, y) to move it
		anchor=north east,  % Attach the legend to its top-right corner
		draw=none,           % Optional: Remove border around the legend
		fill=none,
		align=left
	}
	]
	
	\node[align=center] at (axis cs: 0.267, 2.8) {\footnotesize\sc a};
	%\node[align=center] at (axis cs: 1.105, 3500) {\footnotesize B};
	\node[align=center] at (axis cs: 1.254, 33) {\footnotesize\sc c};
	\node[align=center] at (axis cs: 2.940-0.2, 2) {\footnotesize\sc d};
	\node[align=center] at (axis cs: 3.024, 43) {\footnotesize\sc e};
	\node[align=center] at (axis cs: 3.893, 2.5) {\footnotesize\sc f};
	\node[align=center] at (axis cs: 6.643, 320) {\footnotesize\sc g};
	%\node[align=center] at (axis cs: 6.935, 300) {\footnotesize H};
	
	\addplot[
	color=markrood,
	name path=ypointone,
	] table [col sep = comma, x expr = \thisrowno{0}/1000, y index = 1]{M_y_01.csv};
	
	\addplot[
	color=rainbow2of8,
	name path=yone,
	] table [col sep = comma, x expr = \thisrowno{0}/1000, y index = 1]{M_y_1.csv};
	
	\addplot[
	color=rainbow4of8,
	name path=yfive,
	] table [col sep = comma, x expr = \thisrowno{0}/1000, y index = 1]{M_y_5.csv};
	
	\addplot+[draw=none,domain=0:10,name path=yzero] {0.01}; 
	
	\addplot[color=markrood, opacity=0.5] fill between[
	of = yone and ypointone,
	]; 
	
	\addplot[color=rainbow2of8, opacity=0.5] fill between[
	of = yfive and yone,
	]; 
	
	\addplot[color=rainbow4of8, opacity=0.5] fill between[
	of = yzero and yfive,
	]; 
	
	\node at (rel axis cs: 0.52, 0.65) [anchor=center] {\includegraphics[width=3.5cm]{MARWANFINGERtrans.png}};

	%\addlegendentry{\small \hspace{0.3em} $G=\SI{0.1}{\percent}$}
	%\addlegendentry{\small \hspace{0.3em} $G=\SI{1}{\percent}$}
	%\addlegendentry{\small \hspace{0.3em} $G=\SI{5}{\percent}$}
	
	\draw[thick, <->] (rel axis cs: 0.605, 0.81) --++ (30:0.5cm);
	\node[align=left] at (rel axis cs: 0.615, 0.895) {\footnotesize $y$};
	
	\nextgroupplot
	[
	ymode=log,
	ymin=1e-0,
	ymax=1e4,
	%ymin=1e-2, ymax=1e0,
	xlabel={Enforced acceleration frequency (\si{\kilo\hertz})},
	%ylabel={Max. linear acceleration},
	xmin=0, xmax=10,
	%ymin=-0.1, ymax=1.1,
	%ymin=0, ymax=120,
	%xtick={0,20,40,60,80,100},
	%ytick={0,20,40,60,80,100,120},
	%legend pos=north west,
	ymajorgrids=true,
	xmajorgrids=true,
	grid style=dashed,
	legend cell align={left},
	legend style={
		at={(1, 1)},  % Adjust (x, y) to move it
		anchor=north east,  % Attach the legend to its top-right corner
		draw=none,           % Optional: Remove border around the legend
		fill=none,
		align=left
	}
	]
	
	\node[align=center] at (axis cs: 0.267, 3500) {\footnotesize\sc a};
	%\node[align=center] at (axis cs: 1.105, 3500) {\footnotesize B};
	\node[align=center] at (axis cs: 1.254, 2200) {\footnotesize\sc c};
	\node[align=center] at (axis cs: 2.940-0.2, 30) {\footnotesize\sc d};
	\node[align=center] at (axis cs: 3.024, 1600) {\footnotesize\sc e};
	\node[align=center] at (axis cs: 3.893, 4.5) {\footnotesize\sc f};
	\node[align=center] at (axis cs: 6.643, 1500) {\footnotesize\sc g};
	%\node[align=center] at (axis cs: 6.935, 300) {\footnotesize H};
	
	\addplot[
	color=markrood,
	name path=zpointone,
	] table [col sep = comma, x expr = \thisrowno{0}/1000, y index = 1]{M_z_01.csv};
	
	\addplot[
	color=rainbow2of8,
	name path=zone,
	] table [col sep = comma, x expr = \thisrowno{0}/1000, y index = 1]{M_z_1.csv};
	
	\addplot[
	color=rainbow4of8,
	name path=zfive,
	] table [col sep = comma, x expr = \thisrowno{0}/1000, y index = 1]{M_z_5.csv};
	
	\addplot+[draw=none,domain=0:10,name path=zzero] {0.01}; 
	
	\addplot[color=markrood, opacity=0.5] fill between[
	of = zone and zpointone,
	]; 
	
	\addplot[color=rainbow2of8, opacity=0.5] fill between[
	of = zfive and zone,
	]; 
	
	\addplot[color=rainbow4of8, opacity=0.5] fill between[
	of = zzero and zfive,
	]; 
	
	\node at (rel axis cs: 0.52, 0.65) [anchor=center] {\includegraphics[width=3.5cm]{MARWANFINGERtrans.png}};

	\draw[thick, <->] (rel axis cs: 0.615, 0.79) --++ (90:0.5cm);
	\node[align=left] at (rel axis cs: 0.635, 0.86) {\footnotesize $z$};
	
	
	\end{groupplot}
	
	
	
	\end{tikzpicture}
	\caption[Modal frequency response analysis for fibre-tip cavity mounting structure]{Modal frequency response analysis for fibre-tip cavity mounting structure. The graphs show the maximum linear acceleration magnitude of the structure versus the frequency of enforced acceleration at the base in three different directions: $x$ (top), $y$ (middle) and $z$ (bottom). In each direction, traces are shown for three different values of the structural damping coefficient $G$. The magnitude of the enforced acceleration equals \SI{1}{\milli\meter\per\second\squared}}
	\label{freqresponseMarw} 
\end{figure}



\begin{figure}[t]
	\vspace{0em}
	\centering
	\begin{subfigure}[b]{.45\linewidth}
		\begin{tikzpicture}
		\centering
		\node[inner sep=0] (image) at (0,0) {\includegraphics[width=5cm] {1F.png}};;
		\begin{scope}[x={(image.south east)},y={(image.north west)}]
		\node at (0.3, 0.7) {\footnotesize \SI{267}{\hertz}};
		\end{scope}
		\end{tikzpicture}
		\centering
		\caption{}
		\label{a}
	\end{subfigure}
	\begin{subfigure}[b]{.45\linewidth}
		\begin{tikzpicture}
		\centering
		\node[inner sep=0] (image) at (0,0) {\includegraphics[width=5cm] {2F.png}};;
		\begin{scope}[x={(image.south east)},y={(image.north west)}]
		\node at (0.3, 0.7) {\footnotesize \SI{1105}{\hertz}};
		\end{scope}
		\end{tikzpicture}
		\centering
		\caption{}
		\label{b}
	\end{subfigure}
	\\
	\begin{subfigure}[b]{.45\linewidth}
		\begin{tikzpicture}
		\centering
		\node[inner sep=0] (image) at (0,0) {\includegraphics[width=5cm] {3F.png}};;
		\begin{scope}[x={(image.south east)},y={(image.north west)}]
		\node at (0.3, 0.7) {\footnotesize \SI{1254}{\hertz}};
		\end{scope}
		\end{tikzpicture}
		\centering
		\caption{}
		\label{c}
	\end{subfigure}
	\begin{subfigure}[b]{.45\linewidth}
		\begin{tikzpicture}
		\centering
		\node[inner sep=0] (image) at (0,0) {\includegraphics[width=5cm] {4F.png}};;
		\begin{scope}[x={(image.south east)},y={(image.north west)}]
		\node at (0.3, 0.7) {\footnotesize \SI{2940}{\hertz}};
		\end{scope}
		\end{tikzpicture}
		\centering
		\caption{}
		\label{c}
	\end{subfigure}
	\\
	\begin{subfigure}[b]{.45\linewidth}
		\begin{tikzpicture}
		\centering
		\node[inner sep=0] (image) at (0,0) {\includegraphics[width=5cm] {5F.png}};;
		\begin{scope}[x={(image.south east)},y={(image.north west)}]
		\node at (0.3, 0.7) {\footnotesize \SI{3024}{\hertz}};
		\end{scope}
		\end{tikzpicture}
		\centering
		\caption{}
		\label{a}
	\end{subfigure}
	\begin{subfigure}[b]{.45\linewidth}
		\begin{tikzpicture}
		\centering
		\node[inner sep=0] (image) at (0,0) {\includegraphics[width=5cm] {6F.png}};;
		\begin{scope}[x={(image.south east)},y={(image.north west)}]
		\node at (0.3, 0.7) {\footnotesize \SI{3893}{\hertz}};
		\end{scope}
		\end{tikzpicture}
		\centering
		\caption{}
		\label{b}
	\end{subfigure}
\\
\begin{subfigure}[b]{.45\linewidth}
	\begin{tikzpicture}
	\centering
	\node[inner sep=0] (image) at (0,0) {\includegraphics[width=5cm] {7F.png}};;
	\begin{scope}[x={(image.south east)},y={(image.north west)}]
	\node at (0.3, 0.7) {\footnotesize \SI{6643}{\hertz}};
	\end{scope}
	\end{tikzpicture}
	\centering
	\caption{}
	\label{a}
\end{subfigure}
\begin{subfigure}[b]{.45\linewidth}
	\begin{tikzpicture}
	\centering
	\node[inner sep=0] (image) at (0,0) {\includegraphics[width=5cm] {8F.png}};;
	\begin{scope}[x={(image.south east)},y={(image.north west)}]
	\node at (0.3, 0.7) {\footnotesize \SI{6935}{\hertz}};
		\draw[thick, ->] (1, 0.15) --++ (90:0.5cm) node[above, anchor=south] {\footnotesize $z$};
\draw[thick, ->] (1, 0.15) --++ (150:0.5cm) node[above, anchor=base] {\footnotesize  $x$\,\, \phantom{x}};
\draw[thick, ->] (1, 0.15) --++ (210:0.5cm) node[above, anchor=mid] {\footnotesize $y$\,\, \phantom{x}};
	\end{scope}
	\end{tikzpicture}
	\centering
	\caption{}
	\label{b}
\end{subfigure}
	\caption[Normal modes below \SI{10}{\kilo\hertz} for fibre-tip cavity mounting structure]{Normal modes and corresponding frequencies below \SI{10}{\kilo\hertz} for the fibre-tip cavity mounting structure. The colour scale corresponds to the displacement from the unperturbed position of the structure (dashed contours). Mode labels correspond to those in \autoref{freqresponseMarw}.}
	\label{marwanmodes}
\end{figure}

\begin{table}[h]
	\centering
	\begin{tabular}{ccccccc}
		\toprule
		{Mode number} & { ${T}_x$} & {${T}_y$} & {${T}_z$} & {${R}_x$} & {${R}_y$} & {${R}_z$} \\
		\midrule
		1 (\textsc{a}) & 0.0000 & 0.0000 & 16.5379 & 85.1071 & 0.0000 & 0.0000 \\
		2 (\textsc{b}) & 15.8373 & 0.0000 & 0.0000 & 0.0000 & 1.8747 & 83.4516 \\
		3 (\textsc{c}) & 0.0001 & 0.0014 & 4.9304 & 4.4790 & 0.0003 & 0.0000 \\
		4 (\textsc{d}) & 1.6181 & 0.0004 & 0.0004 & 0.0002 & 0.5771 & 2.7442 \\
		5 (\textsc{e}) & 0.0001 & 0.0024 & 2.7994 & 1.4686 & 2.7770 & 8.9820 \\
		6 (\textsc{f}) & 2.5962 & 0.0000 & 0.0000 & 0.0322 & 1.1527 & 2.7770 \\
		7 (\textsc{g}) & 0.0000 & 0.1034 & 1.5245 & 0.4678 & 0.0199 & 0.0000 \\
		8 (\textsc{h}) & 1.7830 & 0.0000 & 0.0000 & 0.2995 & 0.0134 & 1.1527 \\
		9 & 3.0222 & 0.0037 & 0.5778 & 0.0001 & 0.7470 & 1.3507 \\
		10 & 0.0005 & 8.4142 & 0.4302 & 0.4197 & 0.0001 & 0.0002 \\
		\midrule
		{Total} & 24.8574 & 8.5251 & 26.2227 & 91.9425 & 5.9831 & 91.5980 \\
		\bottomrule
	\end{tabular}
\caption[Percentages of modal effective mass for cavity (fibre-tip) mounting structure]{Modal effective mass participation for the first 10 modes of the cavity (fibre-tip) mounting structure. The numbers shown represent percentages for the translational (${T}_x$, ${T}_y$, ${T}_z$) and rotational (${R}_x$, ${R}_y$, ${R}_z$) degrees of freedom. Labels are references to the modes as indicated in \autoref{freqresponseMarw} and \autoref{marwanmodes}.}
	\label{modeparticipationM}
\end{table}

\begin{figure}[!t]
	%\vspace{0.7em}
	\centering
	\begin{tikzpicture}
	

	
	\begin{groupplot}[group style={group size=1 by 1, vertical sep=1cm, y descriptions at=edge right,},xmin=0,ymin=0,height=5.5cm,width=12cm,no markers,]
	
	\nextgroupplot
	[
	ymode=log,
	ymin=7e-1,
	ymax=1e4,
	ylabel style=markgrijs,
	yticklabel style=markgrijs,
	%ymin=1e-2, ymax=1e0,
	xlabel={Frequency (\si{\kilo\hertz})},
	ylabel={Max. accel. (\si{\milli\meter\per\second\squared})},
	xmin=0, xmax=10,
	ytick pos=right,
	%ymin=-0.1, ymax=1.1,
	%ymin=0, ymax=120,
	%xtick={0,20,40,60,80,100},
	%ytick={0,20,40,60,80,100,120},
	%legend pos=north west,
	ymajorgrids=true,
	xmajorgrids=false,
	grid style=dashed,
	legend cell align={left},
	legend style={
		at={(1, 1)},  % Adjust (x, y) to move it
		anchor=north east,  % Attach the legend to its top-right corner
		draw=none,           % Optional: Remove border around the legend
		fill=none,
		align=left
	}
	]
	
	\draw[color=markgrijs, dashed] 
	(axis cs:0.267, 7e-1) -- (axis cs:0.267, 40000);
	\draw[color=markgrijs, dashed] 
	(axis cs:1.105, 7e-1) -- (axis cs:1.105, 40000);
	\draw[color=markgrijs, dashed] 
	(axis cs:1.254, 7e-1) -- (axis cs:1.254, 40000);
	\draw[color=markgrijs, dashed] 
	(axis cs:2.940, 7e-1) -- (axis cs:2.940, 40000);
	\draw[color=markgrijs, dashed] 
	(axis cs:3.024, 7e-1) -- (axis cs:3.024, 40000);
	\draw[color=markgrijs, dashed] 
	(axis cs:3.893, 7e-1) -- (axis cs:3.893, 40000);
	\draw[color=markgrijs, dashed] 
	(axis cs:6.643, 7e-1) -- (axis cs:6.643, 40000);
	\draw[color=markgrijs, dashed] 
	(axis cs:6.935, 7e-1) -- (axis cs:6.935, 40000);
	
	
	
	\pgfplotsset{
		after end axis/.append code={
			\node[align=center] at (axis cs: 0.267, 20000) {\footnotesize\sc a};
			\node[align=center] at (axis cs: 1.105-0.06, 20000) {\footnotesize\sc b};
			\node[align=center] at (axis cs: 1.254+0.06, 20000) {\footnotesize\sc c};
			\node[align=center] at (axis cs: 2.940-0.1, 20000) {\footnotesize\sc d};
			\node[align=center] at (axis cs: 3.024+0.1, 20000) {\footnotesize\sc e};
			\node[align=center] at (axis cs: 3.893, 20000) {\footnotesize\sc f};
			\node[align=center] at (axis cs: 6.643, 20000) {\footnotesize\sc g};
			\node[align=center] at (axis cs: 6.935, 20000) {\footnotesize\sc h};
		}
	}
	
	
	
	\addplot[
	color=markgrijs,
	name path=xpointone,
	] table [col sep = comma, x expr = \thisrowno{0}/1000, y index = 1]{M_x_01.csv};
	
	\addplot+[draw=none,domain=0:10,name path=xzero] {1e-11}; 
	
	\addplot[color=markgrijs, opacity=0.1] fill between[
	of = xzero and xpointone,
	]; 
	
	\addplot[
	color=markgrijs,
	name path=ypointone,
	] table [col sep = comma, x expr = \thisrowno{0}/1000, y index = 1]{M_y_01.csv};
	
	\addplot+[draw=none,domain=0:10,name path=yzero] {1e-11}; 
	
	\addplot[color=markgrijs, opacity=0.1] fill between[
	of = yzero and ypointone,
	]; 
	
	\addplot[
	color=markgrijs,
	name path=zpointone,
	] table [col sep = comma, x expr = \thisrowno{0}/1000, y index = 1]{M_z_01.csv};
	
	\addplot+[draw=none,domain=0:10,name path=zzero] {1e-11}; 
	
	\addplot[color=markgrijs, opacity=0.1] fill between[
	of = zzero and zpointone,
	]; 
	
	\end{groupplot}
	
	\begin{groupplot}[group style={group size=1 by 1, vertical sep=1cm,y descriptions at=edge left},xmin=0,ymin=0,height=5.5cm,width=12cm,no markers]
	
	\nextgroupplot
	[
	ymode=log,
	ymin=1e-10,
	ymax=1e-4,
	ylabel style=markrood,
	yticklabel style=markrood,
	every y tick/.style={markrood},
	%ymin=1e-2, ymax=1e0,
	%xlabel={Frequency (\si{\kilo\hertz})},
	ylabel={Power spectral density (\si{\volt\squared\per\hertz})},
	xmin=0, xmax=10,
	%ymin=-0.1, ymax=1.1,
	%ymin=0, ymax=120,
	%xtick={0,20,40,60,80,100},
	%ytick={0,20,40,60,80,100,120},
	%legend pos=north west,
	xtick=\empty,
	axis line style=transparent,
	ytick pos=left,
	ymajorgrids=false,
	xmajorgrids=false,
	grid style=dashed,
	y grid style={color=markrood, opacity=0.3},
	legend cell align={left},
	legend style={
		at={(1, 1)},  % Adjust (x, y) to move it
		anchor=north east,  % Attach the legend to its top-right corner
		draw=none,           % Optional: Remove border around the legend
		fill=none,
		align=left
	}
	]
	
	
	\addplot[
	color=markrood,
	] table [col sep = comma, x expr = \thisrowno{0}/1000, y index = 1]{PSD.csv};
	
	
	\end{groupplot}
	
	\end{tikzpicture}
	\caption[Power spectral density of the fibre-tip cavity transmission signal at \SI{1064}{\nano\meter}]{Power spectral density of the fibre-tip cavity transmission signal at \SI{1064}{\nano\meter}. The grey reference traces correspond to the simulated modal frequency response curves ($x$, $y$ and $z$) as shown in \autoref{freqresponseMarw}. Labels are references to the modes as indicated in \autoref{freqresponseMarw} and \autoref{marwanmodes}.} % The text in the square bracket is the caption for the list of figures while the text in the curly brackets is the figure caption
	\label{PSD} 
\end{figure}

\begin{figure}[!t]
	%\vspace{0.7em}
	\centering
	\begin{tikzpicture}
	
	
	
	\begin{groupplot}[group style={group size=1 by 2, vertical sep=1cm, y descriptions at=edge right,},xmin=0,ymin=0,height=5.5cm,width=12cm,no markers,]
	
	\nextgroupplot
	[
	ymode=log,
	ymin=7e-1,
	ymax=1e4,
	ylabel style=markgrijs,
	yticklabel style=markgrijs,
	%ymin=1e-2, ymax=1e0,
	%xlabel={Frequency (\si{\kilo\hertz})},
	ylabel={Max. accel. (\si{\milli\meter\per\second\squared})},
	xmin=0, xmax=10,
	ytick pos=right,
	xticklabels={,,},
	%ymin=-0.1, ymax=1.1,
	%ymin=0, ymax=120,
	%xtick={0,20,40,60,80,100},
	%ytick={0,20,40,60,80,100,120},
	%legend pos=north west,
	ymajorgrids=true,
	xmajorgrids=false,
	grid style=dashed,
	legend cell align={left},
	legend style={
		at={(1, 1)},  % Adjust (x, y) to move it
		anchor=north east,  % Attach the legend to its top-right corner
		draw=none,           % Optional: Remove border around the legend
		fill=none,
		align=left
	}
	]
	
	\draw[color=markgrijs, dashed] 
	(axis cs:0.267, 7e-1) -- (axis cs:0.267, 40000);
	\draw[color=markgrijs, dashed] 
	(axis cs:1.105, 7e-1) -- (axis cs:1.105, 40000);
	\draw[color=markgrijs, dashed] 
	(axis cs:1.254, 7e-1) -- (axis cs:1.254, 40000);
	\draw[color=markgrijs, dashed] 
	(axis cs:2.940, 7e-1) -- (axis cs:2.940, 40000);
	\draw[color=markgrijs, dashed] 
	(axis cs:3.024, 7e-1) -- (axis cs:3.024, 40000);
	\draw[color=markgrijs, dashed] 
	(axis cs:3.893, 7e-1) -- (axis cs:3.893, 40000);
	\draw[color=markgrijs, dashed] 
	(axis cs:6.643, 7e-1) -- (axis cs:6.643, 40000);
	\draw[color=markgrijs, dashed] 
	(axis cs:6.935, 7e-1) -- (axis cs:6.935, 40000);
	
	
	
	\pgfplotsset{
		after end axis/.append code={
			\node[align=center] at (axis cs: 0.267, 20000) {\footnotesize\sc a};
			\node[align=center] at (axis cs: 1.105-0.06, 20000) {\footnotesize\sc b};
			\node[align=center] at (axis cs: 1.254+0.06, 20000) {\footnotesize\sc c};
			\node[align=center] at (axis cs: 2.940-0.1, 20000) {\footnotesize\sc d};
			\node[align=center] at (axis cs: 3.024+0.1, 20000) {\footnotesize\sc e};
			\node[align=center] at (axis cs: 3.893, 20000) {\footnotesize\sc f};
			\node[align=center] at (axis cs: 6.643, 20000) {\footnotesize\sc g};
			\node[align=center] at (axis cs: 6.935, 20000) {\footnotesize\sc h};
		}
	}
	
	
	
	\addplot[
	color=markgrijs,
	name path=xpointone,
	] table [col sep = comma, x expr = \thisrowno{0}/1000, y index = 1]{M_x_01.csv};
	
	\addplot+[draw=none,domain=0:10,name path=xzero] {1e-11}; 
	
	\addplot[color=markgrijs, opacity=0.1] fill between[
	of = xzero and xpointone,
	]; 
	
	\addplot[
	color=markgrijs,
	name path=ypointone,
	] table [col sep = comma, x expr = \thisrowno{0}/1000, y index = 1]{M_y_01.csv};
	
	\addplot+[draw=none,domain=0:10,name path=yzero] {1e-11}; 
	
	\addplot[color=markgrijs, opacity=0.1] fill between[
	of = yzero and ypointone,
	]; 
	
	\addplot[
	color=markgrijs,
	name path=zpointone,
	] table [col sep = comma, x expr = \thisrowno{0}/1000, y index = 1]{M_z_01.csv};
	
	\addplot+[draw=none,domain=0:10,name path=zzero] {1e-11}; 
	
	\addplot[color=markgrijs, opacity=0.1] fill between[
	of = zzero and zpointone,
	]; 
	
	\nextgroupplot
	[
	ymode=log,
	ymin=7e-1,
	ymax=1e4,
	ylabel style=markgrijs,
	yticklabel style=markgrijs,
	%ymin=1e-2, ymax=1e0,
	xlabel={Frequency (\si{\kilo\hertz})},
	ylabel={Max. accel. (\si{\milli\meter\per\second\squared})},
	xmin=0, xmax=10,
	ytick pos=right,
	%ymin=-0.1, ymax=1.1,
	%ymin=0, ymax=120,
	%xtick={0,20,40,60,80,100},
	%ytick={0,20,40,60,80,100,120},
	%legend pos=north west,
	ymajorgrids=true,
	xmajorgrids=false,
	grid style=dashed,
	legend cell align={left},
	legend style={
		at={(1, 1)},  % Adjust (x, y) to move it
		anchor=north east,  % Attach the legend to its top-right corner
		draw=none,           % Optional: Remove border around the legend
		fill=none,
		align=left
	}
	]
	
	\draw[color=markgrijs, dashed] 
	(axis cs:0.267, 7e-1) -- (axis cs:0.267, 40000);
	\draw[color=markgrijs, dashed] 
	(axis cs:1.105, 7e-1) -- (axis cs:1.105, 40000);
	\draw[color=markgrijs, dashed] 
	(axis cs:1.254, 7e-1) -- (axis cs:1.254, 40000);
	\draw[color=markgrijs, dashed] 
	(axis cs:2.940, 7e-1) -- (axis cs:2.940, 40000);
	\draw[color=markgrijs, dashed] 
	(axis cs:3.024, 7e-1) -- (axis cs:3.024, 40000);
	\draw[color=markgrijs, dashed] 
	(axis cs:3.893, 7e-1) -- (axis cs:3.893, 40000);
	\draw[color=markgrijs, dashed] 
	(axis cs:6.643, 7e-1) -- (axis cs:6.643, 40000);
	\draw[color=markgrijs, dashed] 
	(axis cs:6.935, 7e-1) -- (axis cs:6.935, 40000);
	
	
	
	\pgfplotsset{
		after end axis/.append code={
			\node[align=center] at (axis cs: 0.267, 20000) {\footnotesize\sc a};
			\node[align=center] at (axis cs: 1.105-0.06, 20000) {\footnotesize\sc b};
			\node[align=center] at (axis cs: 1.254+0.06, 20000) {\footnotesize\sc c};
			\node[align=center] at (axis cs: 2.940-0.1, 20000) {\footnotesize\sc d};
			\node[align=center] at (axis cs: 3.024+0.1, 20000) {\footnotesize\sc e};
			\node[align=center] at (axis cs: 3.893, 20000) {\footnotesize\sc f};
			\node[align=center] at (axis cs: 6.643, 20000) {\footnotesize\sc g};
			\node[align=center] at (axis cs: 6.935, 20000) {\footnotesize\sc h};
		}
	}
	
	
	
	\addplot[
	color=markgrijs,
	name path=xpointone,
	] table [col sep = comma, x expr = \thisrowno{0}/1000, y index = 1]{M_x_01.csv};
	
	\addplot+[draw=none,domain=0:10,name path=xzero] {1e-11}; 
	
	\addplot[color=markgrijs, opacity=0.1] fill between[
	of = xzero and xpointone,
	]; 
	
	\addplot[
	color=markgrijs,
	name path=ypointone,
	] table [col sep = comma, x expr = \thisrowno{0}/1000, y index = 1]{M_y_01.csv};
	
	\addplot+[draw=none,domain=0:10,name path=yzero] {1e-11}; 
	
	\addplot[color=markgrijs, opacity=0.1] fill between[
	of = yzero and ypointone,
	]; 
	
	\addplot[
	color=markgrijs,
	name path=zpointone,
	] table [col sep = comma, x expr = \thisrowno{0}/1000, y index = 1]{M_z_01.csv};
	
	\addplot+[draw=none,domain=0:10,name path=zzero] {1e-11}; 
	
	\addplot[color=markgrijs, opacity=0.1] fill between[
	of = zzero and zpointone,
	]; 
	
	\end{groupplot}
	
	\begin{groupplot}[group style={group size=1 by 2, vertical sep=1cm,y descriptions at=edge left},xmin=0,ymin=0,height=5.5cm,width=12cm,no markers]
	
	\nextgroupplot
	[
	%ymode=log,
	%ymin=1e-10,
	%ymax=1e-4,
	ylabel style=markrood,
	yticklabel style=markrood,
	every y tick/.style={markrood},
	%ymin=1e-2, ymax=1e0,
	%xlabel={Frequency (\si{\kilo\hertz})},
	ylabel={$\vert H(i\omega)\vert$},
	xmin=0, xmax=10,
	ytick={0,0.05,0.1},
	yticklabels = {0,0.05,0.1},
	%ymin=-0.1, ymax=1.1,
	%ymin=0, ymax=120,
	%xtick={0,20,40,60,80,100},
	%ytick={0,20,40,60,80,100,120},
	%legend pos=north west,
	xtick=\empty,
	axis line style=transparent,
	ytick pos=left,
	ymajorgrids=false,
	xmajorgrids=false,
	grid style=dashed,
	y grid style={color=markrood, opacity=0.3},
	legend cell align={left},
	legend style={
		at={(1, 1)},  % Adjust (x, y) to move it
		anchor=north east,  % Attach the legend to its top-right corner
		draw=none,           % Optional: Remove border around the legend
		fill=none,
		align=left
	}
	]
	
	
	\addplot[
	color=markrood,
	] table [col sep = comma, x expr = \thisrowno{0}/1000, y index = 1]{TRANSFERFUNCTION.csv};
	
	\nextgroupplot
	[
	%ymode=log,
	%ymin=1e-10,
	%ymax=1e-4,
	ylabel style=markdiepblauw,
	yticklabel style=markdiepblauw,
	every y tick/.style={markdiepblauw},
	%ymin=1e-2, ymax=1e0,
	%xlabel={Frequency (\si{\kilo\hertz})},
	ylabel={$\arg{(H(i\omega))}$ (\si{\degree})},
	xmin=0, xmax=10,
	%ymin=-0.1, ymax=1.1,
	%ymin=0, ymax=120,
	%xtick={0,20,40,60,80,100},
	%ytick={0,20,40,60,80,100,120},
	%legend pos=north west,
	xtick=\empty,
	axis line style=transparent,
	ytick pos=left,
	ymajorgrids=false,
	xmajorgrids=false,
	grid style=dashed,
	y grid style={color=markrood, opacity=0.3},
	legend cell align={left},
	legend style={
		at={(1, 1)},  % Adjust (x, y) to move it
		anchor=north east,  % Attach the legend to its top-right corner
		draw=none,           % Optional: Remove border around the legend
		fill=none,
		align=left
	}
	]
	
	
	\addplot[
	color=markdiepblauw,
	] table [col sep = comma, x expr = \thisrowno{0}/1000, y index = 2]{TRANSFERFUNCTION.csv};
	
	\end{groupplot}
	
	\end{tikzpicture}
\caption[Magnitude and phase of the transfer function corresponding to the transmission at \SI{1064}{\nano\meter} of a piezo-driven fibre-tip cavity]{Magnitude (top) and phase (bottom) of the transfer function $H(i\omega)$ corresponding to the transmission signal at \SI{1064}{\nano\meter} of a fibre-tip cavity that is driven sinusoidally with a single sheer piezo. The grey reference traces correspond to the simulated modal frequency response curves ($x$, $y$ and $z$) as shown in \autoref{freqresponseMarw}. Labels are references to the modes as indicated in \autoref{freqresponseMarw} and \autoref{marwanmodes}.}
\label{transferfunction} 
\end{figure}


\begin{figure}[!t]
	%\vspace{0.7em}
	\centering
	\begin{tikzpicture}
	
	\begin{groupplot}[group style={group size=1 by 3, vertical sep=1cm},xmin=0,ymin=0,height=5.5cm,width=13.5cm,no markers]
	
	\iffalse
	\nextgroupplot
	[
	ylabel={},
	xmin=0, xmax=1,
	ymin=0, ymax=1,
	xticklabels={,\textcolor{white}{0},},
	yticklabels={,,},
	%legend pos=north west,
	ymajorgrids=true,
	xmajorgrids=true,
	grid style=dashed,
	%axis lines=none,        % No axis lines
	%tick style=none,        % No ticks
	%xtick=\empty,           % No x-axis ticks
	%ytick=\empty,           % No y-axis ticks
	xlabel={},              % No x-axis label
	ylabel={},              % No y-axis label
	%enlargelimits=true,     % Keeps the plot within the bounds
	]
	
	\addplot graphics [
	xmin=0, xmax=1,  % Stretch to fit the x-axis limits
	ymin=0, ymax=1    % Stretch to fit the y-axis limits
	] {FINALFINGER.png};
	\fi
	
	
	\nextgroupplot
	[
	ymode=log,
	ymin=1e-0,
	ymax=1e4,
	%xlabel={Frequency},
	%ylabel={Max. linear acceleration},
	xmin=0, xmax=10,
	xticklabels={,,},
	%ymin=-0.1, ymax=1.1,
	%ymin=0, ymax=120,
	%xtick={0,20,40,60,80,100},
	%ytick={0,20,40,60,80,100,120},
	%legend pos=north west,
	ymajorgrids=true,
	xmajorgrids=true,
	grid style=dashed,
	legend cell align={left},
	legend style={
		at={(1, 1)},  % Adjust (x, y) to move it
		anchor=north east,  % Attach the legend to its top-right corner
		draw=none,           % Optional: Remove border around the legend
		fill=none,
		align=left
	}
	]
	
	%\node[align=center] at (axis cs: 1.890, 3500) {\footnotesize A};
	\node[align=center] at (axis cs: 1.979, 3500) {\footnotesize\sc b};
	\node[align=center] at (axis cs: 5.893, 2500) {\footnotesize\sc c};
	\node[align=center] at (axis cs: 6.337, 300) {\footnotesize\sc d};
	\node[align=center] at (axis cs: 6.466, 700) {\footnotesize\sc e};
	\node[align=center] at (axis cs: 8.530+0.2, 90) {\footnotesize\sc f};
	
	
	\addplot[
	color=markrood,
	name path=xpointone,
	] table [col sep = comma, x expr = \thisrowno{0}/1000, y index = 1]{F_x_01.csv};
	
	\addplot[
	color=rainbow2of8,
	name path=xone,
	] table [col sep = comma, x expr = \thisrowno{0}/1000, y index = 1]{F_x_1.csv};
	
	\addplot[
	color=rainbow4of8,
	name path=xfive,
	] table [col sep = comma, x expr = \thisrowno{0}/1000, y index = 1]{F_x_5.csv};
	
	\addplot+[draw=none,domain=0:10,name path=xzero] {0.01}; 
	
	\addplot[color=markrood, opacity=0.5] fill between[
	of = xone and xpointone,
	]; 
	
	\addplot[color=rainbow2of8, opacity=0.5] fill between[
	of = xfive and xone,
	]; 
	
	\addplot[color=rainbow4of8, opacity=0.5] fill between[
	of = xzero and xfive,
	]; 
	
	\node at (rel axis cs: 0.37, 0.65) [anchor=center] {\includegraphics[width=3.5cm]{FINALFINGERtrans.png}};
	
	\addlegendentry{\footnotesize \hspace{0.3em} $G=\SI{0.1}{\percent}$}
	\addlegendentry{\footnotesize \hspace{0.3em} $G=\SI{1}{\percent}$}
	\addlegendentry{\footnotesize \hspace{0.3em} $G=\SI{5}{\percent}$}
	
	\draw[thick, <->] (rel axis cs: 0.54, 0.8) --++ (150:0.5cm);
	\node[align=left] at (rel axis cs: 0.532, 0.88) {\footnotesize $x$};
	
	\nextgroupplot
	[
	ymode=log,
	ymin=1e-0,
	ymax=1e2,
	%ymin=1e-2, ymax=1e0,
	%xlabel={Frequency},
	ylabel={Maximum linear acceleration magnitude (\si{\milli\meter\per\second\squared})},
	xmin=0, xmax=10,
	xticklabels={,,},
	%ymin=-0.1, ymax=1.1,
	%ymin=0, ymax=120,
	%xtick={0,20,40,60,80,100},
	%ytick={0,20,40,60,80,100,120},
	%legend pos=north west,
	ymajorgrids=true,
	xmajorgrids=true,
	grid style=dashed,
	legend cell align={left},
	legend style={
		at={(1, 1)},  % Adjust (x, y) to move it
		anchor=north east,  % Attach the legend to its top-right corner
		draw=none,           % Optional: Remove border around the legend
		fill=none,
		align=left
	}
	]
	
	%\node[align=center] at (axis cs: 1.890, 3500) {\footnotesize A};
	\node[align=center] at (axis cs: 1.979, 6) {\footnotesize\sc b};
	\node[align=center] at (axis cs: 5.893, 20) {\footnotesize\sc c};
	\node[align=center] at (axis cs: 6.337, 61) {\footnotesize\sc d};
	\node[align=center] at (axis cs: 6.466, 42) {\footnotesize\sc e};
	\node[align=center] at (axis cs: 8.530, 13) {\footnotesize\sc f};
	
	\addplot[
	color=markrood,
	name path=ypointone,
	] table [col sep = comma, x expr = \thisrowno{0}/1000, y index = 1]{F_y_01.csv};
	
	\addplot[
	color=rainbow2of8,
	name path=yone,
	] table [col sep = comma, x expr = \thisrowno{0}/1000, y index = 1]{F_y_1.csv};
	
	\addplot[
	color=rainbow4of8,
	name path=yfive,
	] table [col sep = comma, x expr = \thisrowno{0}/1000, y index = 1]{F_y_5.csv};
	
	\addplot+[draw=none,domain=0:10,name path=yzero] {0.01}; 
	
	\addplot[color=markrood, opacity=0.5] fill between[
	of = yone and ypointone,
	]; 
	
	\addplot[color=rainbow2of8, opacity=0.5] fill between[
	of = yfive and yone,
	]; 
	
	\addplot[color=rainbow4of8, opacity=0.5] fill between[
	of = yzero and yfive,
	]; 
	
	\node at (rel axis cs: 0.37, 0.65) [anchor=center] {\includegraphics[width=3.5cm]{FINALFINGERtrans.png}};
	
	%\addlegendentry{\small \hspace{0.3em} $G=\SI{0.1}{\percent}$}
	%\addlegendentry{\small \hspace{0.3em} $G=\SI{1}{\percent}$}
	%\addlegendentry{\small \hspace{0.3em} $G=\SI{5}{\percent}$}
	
	\draw[thick, <->] (rel axis cs: 0.505, 0.81) --++ (30:0.5cm);
	\node[align=left] at (rel axis cs: 0.515, 0.895) {\footnotesize $y$};
	
	\nextgroupplot
	[
	ymode=log,
	ymin=1e-0,
	ymax=1e4,
	%ymin=1e-2, ymax=1e0,
	xlabel={Enforced acceleration frequency (\si{\kilo\hertz})},
	%ylabel={Max. linear acceleration},
	xmin=0, xmax=10,
	%ymin=-0.1, ymax=1.1,
	%ymin=0, ymax=120,
	%xtick={0,20,40,60,80,100},
	%ytick={0,20,40,60,80,100,120},
	%legend pos=north west,
	ymajorgrids=true,
	xmajorgrids=true,
	grid style=dashed,
	legend cell align={left},
	legend style={
		at={(1, 1)},  % Adjust (x, y) to move it
		anchor=north east,  % Attach the legend to its top-right corner
		draw=none,           % Optional: Remove border around the legend
		fill=none,
		align=left
	}
	]
	
	\node[align=center] at (axis cs: 1.890, 3500) {\footnotesize\sc a};
	%\node[align=center] at (axis cs: 1.979, 3500) {\footnotesize B};
	\node[align=center] at (axis cs: 5.893, 55) {\footnotesize\sc c};
	\node[align=center] at (axis cs: 6.337, 5000) {\footnotesize\sc d};
	\node[align=center] at (axis cs: 6.466, 2500) {\footnotesize\sc e};
	\node[align=center] at (axis cs: 8.530, 300) {\footnotesize\sc f};
	
	\addplot[
	color=markrood,
	name path=zpointone,
	] table [col sep = comma, x expr = \thisrowno{0}/1000, y index = 1]{F_z_01.csv};
	
	\addplot[
	color=rainbow2of8,
	name path=zone,
	] table [col sep = comma, x expr = \thisrowno{0}/1000, y index = 1]{F_z_1.csv};
	
	\addplot[
	color=rainbow4of8,
	name path=zfive,
	] table [col sep = comma, x expr = \thisrowno{0}/1000, y index = 1]{F_z_5.csv};
	
	\addplot+[draw=none,domain=0:10,name path=zzero] {0.01}; 
	
	\addplot[color=markrood, opacity=0.5] fill between[
	of = zone and zpointone,
	]; 
	
	\addplot[color=rainbow2of8, opacity=0.5] fill between[
	of = zfive and zone,
	]; 
	
	\addplot[color=rainbow4of8, opacity=0.5] fill between[
	of = zzero and zfive,
	]; 
	
	\node at (rel axis cs: 0.37, 0.65) [anchor=center] {\includegraphics[width=3.5cm]{FINALFINGERtrans.png}};
	
	\draw[thick, <->] (rel axis cs: 0.515, 0.79) --++ (90:0.5cm);
	\node[align=left] at (rel axis cs: 0.535, 0.86) {\footnotesize $z$};
	
	%\addlegendentry{\small \hspace{0.3em} $G=\SI{0.1}{\percent}$}
	%\addlegendentry{\small \hspace{0.3em} $G=\SI{1}{\percent}$}
	%\addlegendentry{\small \hspace{0.3em} $G=\SI{5}{\percent}$}
	
	\end{groupplot}
	
	
	
	\end{tikzpicture}
	\caption[Modal frequency response analysis for glass-substrate cavity mounting structure]{Modal frequency response analysis for glass-substrate cavity mounting structure. The graphs show the maximum linear acceleration magnitude of the structure versus the frequency of enforced acceleration at the base in three different directions: $x$ (top), $y$ (middle) and $z$ (bottom). In each direction, traces are shown for three different values of the structural damping coefficient $G$. The magnitude of the enforced acceleration equals \SI{1}{\milli\meter\per\second\squared}}
	\label{freqresponseFF} 
\end{figure}



\begin{figure}[t]
	\vspace{0em}
	\centering
	\begin{subfigure}[b]{.45\linewidth}
		\begin{tikzpicture}
		\centering
		\node[inner sep=0] (image) at (0,0) {\includegraphics[width=5cm] {F1F.png}};;
		\begin{scope}[x={(image.south east)},y={(image.north west)}]
		\node at (0.2, 0.7) {\footnotesize \SI{1890}{\hertz}};
		\end{scope}
		\end{tikzpicture}
		\centering
		\caption{}
		\label{a}
	\end{subfigure}
	\begin{subfigure}[b]{.45\linewidth}
		\begin{tikzpicture}
		\centering
		\node[inner sep=0] (image) at (0,0) {\includegraphics[width=5cm] {F2F.png}};;
		\begin{scope}[x={(image.south east)},y={(image.north west)}]
		\node at (0.2, 0.7) {\footnotesize \SI{1979}{\hertz}};
		\end{scope}
		\end{tikzpicture}
		\centering
		\caption{}
		\label{b}
	\end{subfigure}
	\\
	\begin{subfigure}[b]{.45\linewidth}
		\begin{tikzpicture}
		\centering
		\node[inner sep=0] (image) at (0,0) {\includegraphics[width=5cm] {F3F.png}};;
		\begin{scope}[x={(image.south east)},y={(image.north west)}]
		\node at (0.2, 0.7) {\footnotesize \SI{5893}{\hertz}};
		\end{scope}
		\end{tikzpicture}
		\centering
		\caption{}
		\label{c}
	\end{subfigure}
	\begin{subfigure}[b]{.45\linewidth}
		\begin{tikzpicture}
		\centering
		\node[inner sep=0] (image) at (0,0) {\includegraphics[width=5cm] {F4F.png}};;
		\begin{scope}[x={(image.south east)},y={(image.north west)}]
		\node at (0.2, 0.7) {\footnotesize \SI{6337}{\hertz}};
		\end{scope}
		\end{tikzpicture}
		\centering
		\caption{}
		\label{c}
	\end{subfigure}
	\\
	\begin{subfigure}[b]{.45\linewidth}
		\begin{tikzpicture}
		\centering
		\node[inner sep=0] (image) at (0,0) {\includegraphics[width=5cm] {F5F.png}};;
		\begin{scope}[x={(image.south east)},y={(image.north west)}]
		\node at (0.2, 0.7) {\footnotesize \SI{6466}{\hertz}};
		\end{scope}
		\end{tikzpicture}
		\centering
		\caption{}
		\label{a}
	\end{subfigure}
	\begin{subfigure}[b]{.45\linewidth}
		\begin{tikzpicture}
		\centering
		\node[inner sep=0] (image) at (0,0) {\includegraphics[width=5cm] {F6F.png}};;
		\begin{scope}[x={(image.south east)},y={(image.north west)}]
		\node at (0.2, 0.7) {\footnotesize \SI{8530}{\hertz}};
		\draw[thick, ->] (1, 0.15) --++ (90:0.5cm) node[above, anchor=south] {\footnotesize $z$};
\draw[thick, ->] (1, 0.15) --++ (150:0.5cm) node[above, anchor=base] {\footnotesize  $x$\,\, \phantom{x}};
\draw[thick, ->] (1, 0.15) --++ (210:0.5cm) node[above, anchor=mid] {\footnotesize $y$\,\, \phantom{x}};
		\end{scope}
		\end{tikzpicture}
		\centering
		\caption{}
		\label{b}
	\end{subfigure}
	\caption[Normal modes below \SI{10}{\kilo\hertz} for glass-substrate cavity mounting structure]{Normal modes and corresponding frequencies below \SI{10}{\kilo\hertz} for the glass-substrate cavity mounting structure. The colour scale corresponds to the displacement from the unperturbed position of the structure (dashed contours). Mode labels correspond to those in \autoref{freqresponseFF}.}
	\label{FFmodes}
\end{figure}

\begin{table}[h]
	\centering
	\begin{tabular}{ccccccc}
		\toprule
		{Mode number} & { ${T}_x$} & {${T}_y$} & {${T}_z$} & {${R}_x$} & {${R}_y$} & {${R}_z$} \\
		\midrule
		1 (\textsc{a}) & 0.0007 & 0.0000 & 35.8125 & 1.5245 & 0.0004 & 0.0000 \\
		2 (\textsc{b}) & 35.7347 & 0.0001 & 0.0007 & 0.0000 & 0.0006 & 1.5554 \\
		3 (\textsc{c}) & 26.7161 & 0.0023 & 0.0092 & 0.0073 & 0.2986 & 25.4059 \\
		4 (\textsc{d}) & 0.1826 & 0.0070 & 19.8976 & 19.7244 & 8.5420 & 1.1956 \\
		5 (\textsc{e}) & 0.2574 & 0.0023 & 7.2673 & 7.7576 & 15.4124 & 0.2826 \\
		6 (\textsc{f}) & 0.0940 & 0.0001 & 0.0290 & 0.0322 & 1.1623 & 0.1025 \\
		7 & 0.0120 & 1.0386 & 10.8091 & 15.2692 & 0.7660 & 0.0199 \\
		8 & 0.1064 & 45.3301 & 0.2023 & 0.2995 & 0.0543 & 0.1714 \\
		9 & 8.2676 & 0.5778 & 0.0307 & 0.0477 & 10.6172 & 13.3469 \\
		10 & 3.8187 & 0.2649 & 0.0035 & 0.0056 & 26.4222 & 6.3842 \\
		\midrule
		{Total} & 75.1903 & 47.2232 & 74.0620 & 44.6679 & 63.2762 & 47.4645 \\
		\bottomrule
	\end{tabular}
\caption[Percentages of modal effective mass for cavity (glass-substrate) mounting structure]{Modal effective mass participation for the first 10 modes of the cavity (glass-substrate) mounting structure. The numbers shown represent percentages for the translational (${T}_x$, ${T}_y$, ${T}_z$) and rotational (${R}_x$, ${R}_y$, ${R}_z$) degrees of freedom. Labels are references to the modes as indicated in \autoref{freqresponseFF} and \autoref{FFmodes}.}
	\label{tab:modepercentFinalfinger}   
\end{table}







\section{PYRAMID MOT}
\subsection{Introduction}
For the operation of holographic optical tweezers inside a cavity (see \autoref{exp5}), it is not possible to load microtraps from a MOT positioned between the cavity mirrors due to limited optical access for the MOT beams. The idea is to position the MOT cloud in the neighbourhood of the cavity, from which atoms are then transported towards the focal plane of the trapping light inside the cavity. There are various methods of doing this, which include funnelling atoms ejected with an atomic fountain using a combination of blue detuned cavity modes \cite{Puppe2007}, guiding atoms from a static MOT with a single, red detuned beam \cite{Numann2005} or using an optical lattice that can act as an optical conveyor belt \cite{Gallego2018,Schrader2001}. For this experiment, we propose using an atomic fountain to load a large grid of holographic microtraps in the cavity. This will increase the capture probability from the thermally expanding cloud of atoms.



5 - 2 * ((6. * a * t) ** (1./3.))


\begin{figure}[t]
	\hspace{-0cm}
	\begin{tikzpicture}[remember picture]
		\node[anchor=south west,inner sep=0] (image1) at (0,0) {\includegraphics[width=1\textwidth]{{"pyramidinfinger"}.pdf}};
		
		% left four
		\begin{scope}[x={(image1.south east)},y={(image1.north west)}]
			% draw a grid
			%\draw[help lines,xstep=.1,ystep=.1,overlay] (0,0) grid (1,1);
			% draw ticks
			%\foreach \x in {0,1,...,9} { \node [anchor=north,overlay] at (\x/10,0) {0.\x}; }
			%\foreach \y in {0,1,...,9} { \node [anchor=east,overlay] at (0,\y/10) {0.\y}; }
			
			% labels
			%\node[markgrijs, align=center, rotate=0] at (0.0,0.9) {\footnotesize\sc (a)};
			%\node[markgrijs, align=center, rotate=0] at (0.35,0.9) {\footnotesize\sc (b)};
			%\node[markgrijs, align=center, rotate=0] at (0.7,0.9) {\footnotesize\sc (c)};
			
			\node[markgrijs, align=center, rotate=0] at (0.15,-0.05) {\footnotesize\sc (a)};
			\node[markgrijs, align=center, rotate=0] at (0.5,-0.05) {\footnotesize\sc (b)};
			\node[markgrijs, align=center, rotate=0] at (0.85,-0.05) {\footnotesize\sc (c)};
			
		\end{scope}
		
	\end{tikzpicture}
	\caption[Loading an intracavity dipole trap from a pyramid MOT]{Loading an intracavity dipole trap from a pyramid MOT. (\textsc{a}) A cutaway view of the cavity mounting structure, revealing the cavity mirrors and part of the right-angle mirrors that reflect the cooling beams to form the pyramid MOT. At the apex of the pyramid, a central opening allows an atomic beam to emanate towards the cavity mode. (\textsc{b}) The same structure viewed from the back, showing the gold-coated rear surface of the pyramid mirror mount. The location of the cross-section shown in (\textsc{c}) is indicated in light blue. (\textsc{c}) A cross-sectional view illustrating the dipole trap beam (red) and the cooling beams (green), which are incident at small angles relative to the normal. These beams reflect off the gold coating, creating a region around the dipole trap (yellow) where an optical molasses can form, populated by cold atoms from the pyramid MOT.}
	\label{pyramidloading} 
\end{figure}

\subsection{Aperture diffraction}

\begin{figure}[!t]
	\centering
	\begin{tikzpicture}[remember picture]
	\node[anchor=south west,inner sep=0] (image1) at (0,0) {\includegraphics[width=0.75\textwidth]{{"motdiffraction"}.pdf}};
	\begin{scope}[x={(image1.south east)},y={(image1.north west)}]
	% draw a grid
	%\draw[help lines,xstep=.1,ystep=.1,overlay] (0,0) grid (1,1);
	% draw ticks
	%\foreach \x in {0,1,...,9} { \node [anchor=north,overlay] at (\x/10,0) {0.\x}; }
	%\foreach \y in {0,1,...,9} { \node [anchor=east,overlay] at (0,\y/10) {0.\y}; }
	
	\node[markgrijs, align=center, rotate=0] at (0.13,-0.1) {\footnotesize\sc front view}; 
	\node[markgrijs, align=center, rotate=0] at (0.45,-0.1) {\footnotesize\sc side view}; 
	
	\node[white, align=center, rotate=0, anchor=center] at (0.52,0.5) {\footnotesize \SI{780}{\nano\meter}}; 
	\node[white, align=center, rotate=0, anchor=center] at (0.73,0.5) {\footnotesize \SI{780}{\nano\meter}}; 
	
	\draw[markgrijs, ->] (0.848, 0.2) -- (0.848, 0.4) node[left, pos=0.5] {\footnotesize $a$};
	\draw[markgrijs, ->] (0.848, 0.8) -- (0.848, 0.6) node[above, anchor=mid] {\footnotesize };
	
	\draw[markgrijs, dotted] (0.848, 0.2) -- (0.848, 0.05) node[above, anchor=mid] {};
	\draw[markgrijs, <->] (0.848, 0.1) -- (0.983, 0.1) node[below, pos=0.5] {\footnotesize $d$};
	
	\node[markgrijs, align=center, anchor=center] at (0.6245, 0.5) {\footnotesize $=$};
	
	\end{scope}
	
	\end{tikzpicture}
	\caption[asdf]{Schematic of the pyramidal MOT mirror configuration with an aperture at the apex. The front view shows the four reflective faces around the square central aperture. The side view illustrates the incident and reflected \SI{780}{\nano\meter} laser beam. The diffraction of the reflected beam due to the finite aperture size is equivalently described as diffraction from a square obstruction of size $a$ (right). The diffraction pattern is modelled at various distances $d$.}
	\label{motdiffraction} 
\end{figure}

\begin{figure}[t]
	\centering
	\begin{tikzpicture}
	
	\begin{groupplot}[group style={group size=8 by 8,horizontal sep=0.1cm, vertical sep=0.1cm},height=3cm,width=3cm,no markers]
	
	\pgfplotsinvokeforeach{1,...,64} {
		\nextgroupplot[
		ylabel={},
		xmin=0, xmax=1,
		ymin=0, ymax=1,
		xticklabels={,\textcolor{white}{0},},
		yticklabels={,,},
		ymajorgrids=true,
		xmajorgrids=true,
		grid style=dashed,
		axis lines=none,
		xlabel={},
		ylabel={},
		]
		\addplot graphics [
		xmin=0, xmax=1,
		ymin=0, ymax=1
		] {motdiff#1.pdf}; % #1 is the current number
		
		\coordinate (topnode#1) at (axis cs: -0.1,1.1);
		\coordinate (bottomnode#1) at (axis cs: -0.1,-0.15);
		\coordinate (rightnode#1) at (axis cs: 1.15,1.1);
		
		\coordinate (vert#1) at (axis cs: -0.35,0.5);
		\coordinate (hor#1) at (axis cs: 0.5,1.25);
		
		\coordinate (edgeleft#1) at (axis cs: 0.11,0.5);
		\coordinate (edgeright#1) at (axis cs: 0.89,0.5);
		
		\coordinate (fresnel#1) at (axis cs: 0.5,0.85);
		\coordinate (mfresnel#1) at (axis cs: 0.5,0.5);
	}
	

	
	\end{groupplot}
	
	% Set constants
	\pgfmathsetmacro{\lambda}{0.780} % Wavelength in millimetres
	
	% Fresnel numbers for combinations of a and d
	% The pattern: a from {0.1,...,0.8}, d from {0,...,7} → 64 values total
	\pgfplotsinvokeforeach{1,...,64}{
		
		% Convert #1 to indices for a and d:
		% a-index: floor((#1-1)/8), d-index: mod(#1-1,8)
		\pgfmathtruncatemacro{\aindex}{int((#1-1)/8)}
		\pgfmathtruncatemacro{\dindex}{mod(#1-1,8)}
		
		% Map indices to actual values:
		\pgfmathsetmacro{\a}{0.1*(\aindex+1)} % a in {0.1,...,0.8}
		\pgfmathsetmacro{\d}{\dindex}          % d in {0,...,7}
		
		% Calculate Fresnel number:
		\ifnum\dindex=0
			\pgfmathsetmacro{\result}{0}
		\else
			\pgfmathsetmacro{\boolean}{ifthenelse(\a<0.5, 0, 1)} 
			\ifnum\boolean=0
				\pgfmathsetmacro{\result}{1e3*(\a^2)/(\lambda*\d)}
				\node[align=center, white] at (fresnel#1) {\tiny $\pgfmathprintnumber[precision=1]{\result}$};
			\else
				\pgfmathsetmacro{\result}{1e3*(\a^2)/(\lambda*\d)}
				\node[align=center, white] at (mfresnel#1) {\tiny $\pgfmathprintnumber[precision=1]{\result}$};
			\fi
		\fi
	}
	
	
	% Print with three significant digits
	
	

	
	\draw[thick, white, <->] (edgeleft57) -- (edgeright57) node[below, pos=0.5] {\small $a$};
	
	\draw[thick, ->, line cap=rect] (topnode1) -- (bottomnode57) node[left, rotate=0] {$a$};
	\draw[thick, ->, line cap=rect] (topnode1) -- (rightnode8) node[above] {\small $d$};
	
	
	\node[align=center] at (vert1) {\small $0.1$};
	\node[align=center] at (vert9) {\small $0.2$};
	\node[align=center] at (vert17) {\small $0.3$};
	\node[align=center] at (vert25) {\small $0.4$};
	\node[align=center] at (vert33) {\small $0.5$};
	\node[align=center] at (vert41) {\small $0.6$};
	\node[align=center] at (vert49) {\small $0.7$};
	\node[align=center] at (vert57) {\small $0.8$};
	
	\node[align=center] at (hor1) {\small $0$};
	\node[align=center] at (hor2) {\small $1$};
	\node[align=center] at (hor3) {\small $2$};
	\node[align=center] at (hor4) {\small $3$};
	\node[align=center] at (hor5) {\small $4$};
	\node[align=center] at (hor6) {\small $5$};
	\node[align=center] at (hor7) {\small $6$};
	\node[align=center] at (hor8) {\small $7$};
	
	\end{tikzpicture}
	\caption[Simulated diffraction patterns from a pyramid MOT mirror configuration]{Simulated diffraction patterns from a pyramid MOT mirror configuration with a square aperture of various dimensions $a$ at the apex, observed at distances $d$. The Fresnel number $F$ is indicated in each subplot. Each pattern represents a $\SI{1}{\milli\meter}\times\SI{1}{\milli\meter}$ area in the observation plane.}
\label{diffraction2d}
\end{figure}




\begin{figure}[t]
	\centering
	\begin{tikzpicture}
	
	\pgfplotstableread[col sep=comma]{differrs.csv}\datatable
		
		\begin{groupplot}[group style={group size=8 by 8,horizontal sep=0.1cm, vertical sep=0.1cm},height=3cm,width=3cm,no markers]
		
		\pgfplotsinvokeforeach{1,...,64} {
			%\pgfmathtruncatemacro{\rowindex}{1+floor((#1-1)/8)}
			% Plots with indices 8,16,24,...,64 -> show y-tick labels on the right
			\nextgroupplot[
			ylabel={},
			ymin=0, ymax=2,
			xmin=0, xmax=1,
			xticklabels={,\textcolor{white}{0},},
			yticklabels={,,},
			yticklabel pos=right,
			y axis line style={draw=black},
			ymajorgrids=false,
			xmajorgrids=false,
			grid style=dashed,
			xlabel={},
			xtick=\empty,
			ytick=\empty,
			]
			
			
			% Load data from CSV files (motdiffcs#.csv with x,y columns)
			\addplot[
			name path=curve,
			color=markrood,
			] table [x index=0, y index=1, col sep=comma] {motdiffcs#1.csv};
			
			\addplot[
			name path=axis,
			draw=none
			] coordinates {(0,0) (1,0)};
			
			\addplot[
			fill=markrood,
			fill opacity=0.3
			] fill between [of=curve and axis];
			
			\draw[dashed, markgrijs] (0,1) -- (1,1);
			
			\coordinate (topnode#1) at (axis cs:-0.1,2.2);
			\coordinate (bottomnode#1) at (axis cs:-0.1,-0.3);
			\coordinate (rightnode#1) at (axis cs:1.15,2.2);
			
			\coordinate (vert#1) at (axis cs:-0.35,1);
			\coordinate (hor#1) at (axis cs:0.5,2.5);
			
			\coordinate (edgeleft#1) at (axis cs:0.11,0.7);
			\coordinate (edgeright#1) at (axis cs:0.89,0.7);
			
			\coordinate (hedgeleft#1) at (axis cs:0,1.7);
			\coordinate (hedgeright#1) at (axis cs:1,1.7);
			
			\coordinate (outerright#1) at (axis cs:1.15,1);
			\coordinate (louterright#1) at (axis cs:1.15,0);
			
			\pgfmathtruncatemacro{\aindex}{int((#1-1)/8)}
			\pgfmathtruncatemacro{\dindex}{mod(#1-1,8)}
			
			% Map indices to actual values:
			\pgfmathsetmacro{\a}{0.1*(\aindex+1)} % a in {0.1,...,0.8}
			\pgfmathsetmacro{\d}{\dindex}          % d in {0,...,7}
			
			% Calculate Fresnel number:
			\ifnum\dindex=0
				\pgfmathsetmacro{\result}{0}
			\else
				\pgfmathparse{#1 - 1}
				\pgfplotstablegetelem{\pgfmathresult}{[index] 0}\of\datatable % Get nth row, 1st column
				\edef\value{\pgfplotsretval} % Store value
				\pgfmathsetmacro{\result}{\value}
				\node[align=center, markgrijs] at (axis cs: 0.5, 1.7) {\tiny $\pgfmathprintnumber[precision=1]{\result}N^{-1}$};
			\fi
		}
		
		\end{groupplot}
		
		
		
		\draw[thick, <->] (edgeleft57) -- (edgeright57) node[below, pos=0.5] {\footnotesize $a$};
		\draw[thick, <->] (hedgeleft57) -- (hedgeright57) node[below, pos=0.5] {\footnotesize \SI{1}{\milli\meter}};
		
		\draw[thick, ->, line cap=rect] (topnode1) -- (bottomnode57) node[left, rotate=0] {\small $a$};
		\draw[thick, ->, line cap=rect] (topnode1) -- (rightnode8) node[above] {\small $d$};
		
		
		\node[align=center] at (vert1) {\small $0.1$};
		\node[align=center] at (vert9) {\small $0.2$};
		\node[align=center] at (vert17) {\small $0.3$};
		\node[align=center] at (vert25) {\small $0.4$};
		\node[align=center] at (vert33) {\small $0.5$};
		\node[align=center] at (vert41) {\small $0.6$};
		\node[align=center] at (vert49) {\small $0.7$};
		\node[align=center] at (vert57) {\small $0.8$};
		
		\node[align=center] at (hor1) {\small $0$};
		\node[align=center] at (hor2) {\small $1$};
		\node[align=center] at (hor3) {\small $2$};
		\node[align=center] at (hor4) {\small $3$};
		\node[align=center] at (hor5) {\small $4$};
		\node[align=center] at (hor6) {\small $5$};
		\node[align=center] at (hor7) {\small $6$};
		\node[align=center] at (hor8) {\small $7$};
		
		\node[markgrijs, align=center] at (outerright8) {\small $1$};
		\node[markgrijs, align=center] at (outerright16) {\small $1$};
		\node[markgrijs, align=center] at (outerright24) {\small $1$};
		\node[markgrijs, align=center] at (outerright32) {\small $1$};
		\node[markgrijs, align=center] at (outerright40) {\small $1$};
		\node[markgrijs, align=center] at (outerright48) {\small $1$};
		\node[markgrijs, align=center] at (outerright56) {\small $1$};
		\node[markgrijs, align=center] at (outerright64) {\small $1$};
		
		\node[markgrijs, align=center] at (louterright8) {\small $0$};
		\node[markgrijs, align=center] at (louterright16) {\small $0$};
		\node[markgrijs, align=center] at (louterright24) {\small $0$};
		\node[markgrijs, align=center] at (louterright32) {\small $0$};
		\node[markgrijs, align=center] at (louterright40) {\small $0$};
		\node[markgrijs, align=center] at (louterright48) {\small $0$};
		\node[markgrijs, align=center] at (louterright56) {\small $0$};
		\node[markgrijs, align=center] at (louterright64) {\small $0$};
		
	\end{tikzpicture}
	\caption[Central cross-sections of the diffraction profiles from a pyramid MOT mirror configuration]{Central cross-sections of the diffraction profiles from \autoref{diffraction2d}, showing the normalised intensity (such that $1$ corresponds to the incident intensity) along a \SI{1}{\milli\meter} region in front of the aperture for various values of $a$ and $d$.}
\label{diffraction1d}
\end{figure}




\begin{figure}[t]
	\hspace{-0.15cm}
	\begin{tikzpicture}
	\begin{groupplot}[group style={group size=3 by 4,horizontal sep=0.9cm, vertical sep=1.7cm},height=5cm,width=5cm,no markers]
	
		\nextgroupplot
	[
	%xlabel={$x$},
	%ylabel={$y$},
	enlargelimits=false,
	%xmin=0.5,
	%xmax=4.5,
	%ymin=0.5,
	%ymax=4.5,
	axis on top,    % <-- (use with `matrix plot*')
	view={0}{90},
	colormap/viridis,
	colorbar horizontal,
	colorbar style={
		at={(0,1.2)},
		anchor=north west,
		xticklabel pos=upper,
		%x tick scale label style={yshift=-25pt, xshift=23pt},
	},
	point meta=explicit,
	point meta min=0.0,
	point meta max=2.0,
	%xtick={1,2,3,4},
	%ytick={1,2,3,4},
	%xticklabels={1,2,3,4},
	%yticklabels={1,2,3,4},
	]
	\addplot graphics [
	xmin=0, xmax=10,  % Stretch to fit the x-axis limits
	ymin=0, ymax=10    % Stretch to fit the y-axis limits
	] {fresnelfull.pdf};
	
	\draw[fill=none, white] (axis cs:4.5,4.5) rectangle (axis cs:5.5,5.5);
	
	\pgfplotsset{
	after end axis/.append code={
		\node[align=center] at (axis cs: 5, 15) {\footnotesize\sc Fresnel approx.};
		}
	}
	
	
	
	\nextgroupplot
	[
	%xlabel={$x$},
	%ylabel={$y$},
	enlargelimits=false,
	%xmin=0.5,
	%xmax=4.5,
	%ymin=0.5,
	%ymax=4.5,
	axis on top,    % <-- (use with `matrix plot*')
	view={0}{90},
	colormap/viridis,
	colorbar horizontal,
	colorbar style={
		at={(0,1.2)},
		anchor=north west,
		xticklabel pos=upper,
		%x tick scale label style={yshift=-25pt, xshift=23pt},
	},
	point meta=explicit,
	point meta min=0.0,
	point meta max=2.0,
	%xtick={1,2,3,4},
	%ytick={1,2,3,4},
	%xticklabels={1,2,3,4},
	%yticklabels={1,2,3,4},
	]
	\addplot graphics [
	xmin=0, xmax=10,  % Stretch to fit the x-axis limits
	ymin=0, ymax=10    % Stretch to fit the y-axis limits
	] {kirfull.pdf};
	
	\draw[fill=none, white] (axis cs:4.5,4.5) rectangle (axis cs:5.5,5.5);
	
	\pgfplotsset{
	after end axis/.append code={
		\node[align=center] at (axis cs: 5, 15) {\footnotesize\sc direct integration};
		}
	}
	
	
		\nextgroupplot
	[
	%xlabel={$x$},
	%ylabel={$y$},
	enlargelimits=false,
	%xmin=0.5,
	%xmax=4.5,
	%ymin=0.5,
	%ymax=4.5,
	axis on top,    % <-- (use with `matrix plot*')
	view={0}{90},
	colormap/viridis,
	colorbar horizontal,
	colorbar style={
		at={(0,1.2)},
		anchor=north west,
		xticklabel pos=upper,
		xtick={0,2.5,5,7.5,10},
		xticklabels={0,2.5,5,7.5,10}
		%x tick scale label style={yshift=-25pt, xshift=23pt},
	},
	point meta=explicit,
	point meta min=0.0,
	point meta max=10,
	%xtick={1,2,3,4},
	%ytick={1,2,3,4},
	%xticklabels={1,2,3,4},
	%yticklabels={1,2,3,4},
	]
	\addplot graphics [
	xmin=0, xmax=10,  % Stretch to fit the x-axis limits
	ymin=0, ymax=10    % Stretch to fit the y-axis limits
	] {errorfull.pdf};
	
	\draw[fill=none, white] (axis cs:4.5,4.5) rectangle (axis cs:5.5,5.5);
	
	\pgfplotsset{
		after end axis/.append code={
			\node[align=center] at (axis cs: 5, 15) {\footnotesize\sc error (\si{\percent})};
		}
	}
	
	
	
	\nextgroupplot
	[
	%xlabel={$x$},
	ylabel={$y$ (\si{\milli\meter})},
	ylabel style={xshift=2.5cm},
	enlargelimits=false,
	%xmin=0.5,
	%xmax=4.5,
	%ymin=0.5,
	%ymax=4.5,
	axis on top,    % <-- (use with `matrix plot*')
	view={0}{90},
	colormap/viridis,
	colorbar horizontal,
	colorbar style={
		at={(0,1.2)},
		anchor=north west,
		xticklabel pos=upper,
		%x tick scale label style={yshift=-25pt, xshift=23pt},
	},
	point meta=explicit,
	point meta min=0.0,
	point meta max=2.0,
	%xtick={1,2,3,4},
	%ytick={1,2,3,4},
	%xticklabels={1,2,3,4},
	%yticklabels={1,2,3,4},
	]
	\addplot graphics [
	xmin=0, xmax=1,  % Stretch to fit the x-axis limits
	ymin=0, ymax=1    % Stretch to fit the y-axis limits
	] {fresnelzoom.pdf};
	
	\addplot graphics [
	xmin=0.6, xmax=0.98,  % Stretch to fit the x-axis limits
	ymin=0.6, ymax=0.98    % Stretch to fit the y-axis limits
	] {fresnelsmall.pdf};
	
	\draw[fill=none, white] (axis cs:0.6,0.6) rectangle (axis cs:0.98,0.98);
	
	
	\node (lowerleft) at (axis cs: 0, -1.5) {};
	
	
		\nextgroupplot
	[
	xlabel={$x$ (\si{\milli\meter})},
	%ylabel={$y$},
	enlargelimits=false,
	%xmin=0.5,
	%xmax=4.5,
	%ymin=0.5,
	%ymax=4.5,
	axis on top,    % <-- (use with `matrix plot*')
	view={0}{90},
	colormap/viridis,
	colorbar horizontal,
	colorbar style={
		at={(0,1.2)},
		anchor=north west,
		xticklabel pos=upper,
		%x tick scale label style={yshift=-25pt, xshift=23pt},
	},
	point meta=explicit,
	point meta min=0.0,
	point meta max=2.0,
	%xtick={1,2,3,4},
	%ytick={1,2,3,4},
	%xticklabels={1,2,3,4},
	%yticklabels={1,2,3,4},
	]
	\addplot graphics [
	xmin=0, xmax=1,  % Stretch to fit the x-axis limits
	ymin=0, ymax=1    % Stretch to fit the y-axis limits
	] {kirzoom.pdf};

	\addplot graphics [
	xmin=0.6, xmax=0.98,  % Stretch to fit the x-axis limits
	ymin=0.6, ymax=0.98    % Stretch to fit the y-axis limits
	] {kirsmall.pdf};
	
	\draw[fill=none, white] (axis cs:0.6,0.6) rectangle (axis cs:0.98,0.98);
	
	

	
		\nextgroupplot
	[
	%xlabel={$x$ (\si{milli\meter})},
	%ylabel={$y$},
	enlargelimits=false,
	%xmin=0.5,
	%xmax=4.5,
	%ymin=0.5,
	%ymax=4.5,
	axis on top,    % <-- (use with `matrix plot*')
	view={0}{90},
	colormap/viridis,
	colorbar horizontal,
	colorbar style={
		at={(0,1.2)},
		anchor=north west,
		xticklabel pos=upper,
		xtick={0,2.5,5,7.5,10},
		xticklabels={0,2.5,5,7.5,10}
		%x tick scale label style={yshift=-25pt, xshift=23pt},
	},
	point meta=explicit,
	point meta min=0.0,
	point meta max=10,
	%xtick={1,2,3,4},
	%ytick={1,2,3,4},
	%xticklabels={1,2,3,4},
	%yticklabels={1,2,3,4},
	]
	\addplot graphics [
	xmin=0, xmax=1,  % Stretch to fit the x-axis limits
	ymin=0, ymax=1    % Stretch to fit the y-axis limits
	] {errorzoom.pdf};
	
	
	\addplot graphics [
	xmin=0.6, xmax=0.98,  % Stretch to fit the x-axis limits
	ymin=0.6, ymax=0.98    % Stretch to fit the y-axis limits
	] {errorsmallbright.pdf};
	
	\draw[fill=none, white] (axis cs:0.6,0.6) rectangle (axis cs:0.98,0.98);
	
	
	
	\end{groupplot}
	\end{tikzpicture}
	
	\caption[Error analysis of simulated diffraction profiles from a pyramid MOT mirror configuration]{Error analysis of simulated diffraction profiles from a pyramid MOT mirror configuration with an aperture of $\SI{0.1}{\milli\meter}\times\SI{0.1}{\milli\meter}$. The methods are based on the Fresnel approximation (left column) and direct numerical integration of the Fresnel–Kirchhoff integral (middle column). The error between these methods is shown as a percentage of maximum intensity of the undiffracted beam (right column). The top and bottom rows correspond to simulated areas of $\SI{10}{\milli\meter}\times\SI{10}{\milli\meter}$ and a magnification of the central $\SI{1}{\milli\meter}\times\SI{1}{\milli\meter}$ area, respectively. Insets show additional simulations based on the central area only, highlighting significant boundary artefacts in the direct integration approach.}
	\label{differror}
\end{figure}


\begin{figure}[t]
	\vspace{0em}
	\centering
	\begin{tikzpicture}
	\begin{groupplot}[group style={group size=2 by 2,horizontal sep=0.5cm, vertical sep=0.5cm},height=4cm,x=1.05cm,no markers, y label style={at={(-0.085,0.5)}, anchor=base}]
	
	\nextgroupplot
	[
	ymin=0,
	ymax=0.4,
	%xlabel={$x$ (\si{\milli\meter})},
	ylabel={Error (\si{\percent})},
	xticklabels={,,},
	ytick={0,0.1,0.2,0.3,0.4},
	yticklabels={0,10,20,30,40},
	enlargelimits=false,
	ymajorgrids=true,
	xmajorgrids=true,
	grid style=dashed,
	]
	
	\fill[rainbow7of8, opacity=0.2] (axis cs: 4.5, 0) rectangle (axis cs: 5.5, 2);
	
	\addplot[
	name path=curve,
	color=rainbow3of8,
	] table [x index=0, y index=1, col sep=comma] {E.csv};
	
	
	
	\nextgroupplot
	[
	ymin=0,
	ymax=0.4,
	%xlabel={$x$ (\si{\milli\meter})},
	%ylabel={Error (\si{\percent})},
	xtick={0,1},
	xticklabels={,,},
	ytick={0,0.1,0.2,0.3,0.4},
	yticklabels={,,},
	enlargelimits=false,
	ymajorgrids=true,
	xmajorgrids=true,
	grid style=dashed,
	]
	
	\fill[rainbow7of8, opacity=0.2] (axis cs: 0, 0) rectangle (axis cs: 1, 2);
	
	\addplot[
	name path=curve,
	color=rainbow3of8,
	] table [x index=0, y index=1, col sep=comma] {Es.csv};
	
	\nextgroupplot
	[
	ymin=0,
	ymax=2,
	xlabel={$x$ (\si{\milli\meter})},
	ylabel={Norm. intensity},
	enlargelimits=false,
	ymajorgrids=true,
	xmajorgrids=true,
	grid style=dashed,
	legend cell align={left},
	legend style={
		at={(1, 1)},  % Adjust (x, y) to move it
		anchor=north east,  % Attach the legend to its top-right corner
		draw=none,           % Optional: Remove border around the legend
		fill=none
	}
	]
	
	\fill[rainbow7of8, opacity=0.2] (axis cs: 4.5, 0) rectangle (axis cs: 5.5, 2);
	
	\addplot[
	name path=curve,
	color=markdiepblauw,
	opacity=1
	] table [x index=0, y index=1, col sep=comma] {Fcs.csv};
	
	\addplot[
	name path=curve,
	color=markrood,
	opacity=0.7
	] table [x index=0, y index=1, col sep=comma] {Kcs.csv};
	
	
	\legend{\footnotesize\hspace{0.3em}Fresnel approx., \footnotesize\hspace{0.3em}direct integration}
	
	\nextgroupplot
	[
	ymin=0,
	ymax=2,
	%xlabel={$x$ (\si{\milli\meter})},
	%ylabel={Norm. intensity},
	enlargelimits=false,
	ymajorgrids=true,
	xmajorgrids=true,
	grid style=dashed,
	yticklabels={,,},
	xtick={0,1},
	]
	
	\fill[rainbow7of8, opacity=0.2] (axis cs: 0, 0) rectangle (axis cs: 1, 2);
	
	\addplot[
	name path=curve,
	color=markdiepblauw,
	opacity=1
	] table [x index=0, y index=1, col sep=comma] {Fscs.csv};
	
	\addplot[
	name path=curve,
	color=markrood,
	opacity=0.7
	] table [x index=0, y index=1, col sep=comma] {Kscs.csv};
	
	\end{groupplot}
	\end{tikzpicture}
	
	\caption[Cross-sectional error analysis of simulated diffraction profiles from a pyramid MOT mirror configuration]{Cross-sectional error analysis of the diffraction profiles from \autoref{differror} based on the Fresnel approximation and direct numerical integration of the Fresnel–Kirchhoff integral. The error (top) between these methods (bottom) is shown as a percentage of maximum intensity of the undiffracted beam. The left and right columns correspond to simulated areas of $\SI{10}{\milli\meter}\times\SI{10}{\milli\meter}$ and $\SI{1}{\milli\meter}\times\SI{1}{\milli\meter}$ (centred), respectively.}
	\label{differrorcross}
\end{figure}


\subsection{Mirror fabrication}


\begin{figure}[t]
	\centering
	\begin{tikzpicture}[remember picture]
		\node[anchor=south west,inner sep=0] (image1) at (0,0) {\includegraphics[width=0.97\textwidth]{{"polish"}.pdf}};
		
		% left four
		\begin{scope}[x={(image1.south east)},y={(image1.north west)}]
			% draw a grid
			%\draw[help lines,xstep=.1,ystep=.1,overlay] (0,0) grid (1,1);
			% draw ticks
			%\foreach \x in {0,1,...,9} { \node [anchor=north,overlay] at (\x/10,0) {0.\x}; }
			%\foreach \y in {0,1,...,9} { \node [anchor=east,overlay] at (0,\y/10) {0.\y}; }
			
			% labels
			%\node[markgrijs, align=center, rotate=0] at (0.0,0.9) {\footnotesize\sc (a)};
			%\node[markgrijs, align=center, rotate=0] at (0.35,0.9) {\footnotesize\sc (b)};
			%\node[markgrijs, align=center, rotate=0] at (0.7,0.9) {\footnotesize\sc (c)};
			
			\node[markgrijs, align=center, rotate=0] at (0.155,-0.05) {\footnotesize\sc (a)};
			\node[markgrijs, align=center, rotate=0] at (0.5,-0.05) {\footnotesize\sc (b)};
			\node[markgrijs, align=center, rotate=0] at (0.845,-0.05) {\footnotesize\sc (c)};
			
		\end{scope}
		
	\end{tikzpicture}
	\caption[Polishing procedure for right-angle prism mirrors]{Polishing procedure for right-angle prism mirrors. (\textsc{a}) The right-angle prism mirror is secured in a custom aluminium mount using a grub screw. A polishing film is attached to a rotating disk of a fibre optic polishing tool. After removal of the required amount of glass from one corner (\textsc{b}), the prism mirror is flipped inside the mount to polish the adjacent corner (\textsc{c}). White arrows indicate the direction of polishing.}
	\label{polish} 
\end{figure}



\begin{figure}
	\centering
	\begin{tikzpicture}
		\begin{groupplot}[group style={group size=2 by 1,horizontal sep=2.5cm, vertical sep=2.5cm},height=5.5cm,width=5.5cm,no markers]
			
			
			\nextgroupplot
			[
			x=150,
			y=150,
			ylabel={},
			xmin=0, xmax=1,
			ymin=0, ymax=1,
			xticklabels={,\textcolor{white}{0},},
			yticklabels={,,},
			%legend pos=north west,
			ymajorgrids=true,
			xmajorgrids=true,
			grid style=dashed,
			axis lines=none,        % No axis lines
			%tick style=none,        % No ticks
			%xtick=\empty,           % No x-axis ticks
			%ytick=\empty,           % No y-axis ticks
			xlabel={},              % No x-axis label
			ylabel={},              % No y-axis label
			%enlargelimits=true,     % Keeps the plot within the bounds
			clip=false
			]
			
			\addplot graphics [
			xmin=0, xmax=1,  % Stretch to fit the x-axis limits
			ymin=-0.1, ymax=0.9    % Stretch to fit the y-axis limits
			] {mirrorcut.pdf};
			
			\node[markgrijs, align=center, rotate=0] at (0.23,-0.025) {\footnotesize $5$};
			\node[markgrijs, align=center, rotate=0] at (0.413,0.05) {\footnotesize $b$};
			\node[markgrijs, align=center, rotate=0] at (0.255,0.115) {\footnotesize $a$};
			
			\node[markgrijs, align=center, rotate=0] at (0.93,0.52) {\footnotesize $5$};
			\node[markgrijs, align=center, rotate=0] at (0.832,0.175) {\footnotesize $5$};
			
			\node[markrood, align=center, rotate=0] at (0.57,0.23) {\footnotesize $V$};
			
			
			
			
			\nextgroupplot
			[
			ymin=0,
			ymax=5,
			xmin=0,
			xmax=200,
			ytick={0,1,2,3,4,5},
			xlabel={Polishing time $t$ (\si{\second})},
			ylabel={End facet length $a$ (\si{\milli\meter})},
			ymajorgrids=true,
			xmajorgrids=true,
			grid style=dashed,
			name=plot1,
			]
			\addplot[only marks, markgrijs, error bars/.cd,
			y dir=both, y explicit,
			x dir=both, x explicit] table[x=x, y=y, x error=xe, y error=ye] {
				x     y      xe     ye
				168   0.310  0.1    0.005
				168   0.274  0.1    0.005
				168   0.337  0.1    0.005
				168   0.327  0.1    0.005
				165   0.344  0.1    0.005
				169   0.267  0.1    0.005
				180   0.208  0.1    0.005
			};
			
			\addplot[markrood, thick, domain=0:200, samples=200] {(5 - 2 * ( (6 * 0.01280755 * x)^(1/3) ))};
			
			
			
			%\fill[fill=markdiepblauw, opacity=0.5] (axis cs:163,0.190) rectangle (axis cs:182,0.360);
			
			
			
		\end{groupplot}
		\begin{axis}[
			xticklabel style={rotate=90},
			xlabel style={yshift=0.2cm},
			ylabel style={yshift=-0.1cm},
			xmin=163,
			xmax=182,
			ymin=190,
			ymax=360,
			at={(plot1.north east)},
			xshift=-5pt,
			yshift=-5pt,
			width=3.5cm,
			height=3.5cm,
			anchor=north east,
			xlabel={$t$ (\si{\second})},
			ylabel={$a$ (\si{\micro\meter})},
			]
			\fill[fill=white] (axis cs:163,190) rectangle (axis cs:182,360);
			
		\end{axis}
		\begin{axis}[
			xticklabel style={rotate=90},
			xlabel style={yshift=0.2cm},
			ylabel style={yshift=-0.1cm},
			xmin=163,
			xmax=182,
			ymin=190,
			ymax=360,
			at={(plot1.north east)},
			xshift=-5pt,
			yshift=-5pt,
			width=3.5cm,
			height=3.5cm,
			anchor=north east,
			xlabel={$t$ (\si{\second})},
			ylabel={$a$ (\si{\micro\meter})},
			]
			
			\addplot[only marks, markgrijs, error bars/.cd,
			y dir=both, y explicit,
			x dir=both, x explicit] table[x=x, y=y, x error=xe, y error=ye] {
				x     y      xe     ye
				168   310  0.1    5
				168   274  0.1    5
				168   337  0.1    5
				168   327  0.1    5
				165   344  0.1    5
				169   267  0.1    5
				180   208  0.1    5
			};
			
			\addplot[markrood, thick, domain=160:185, samples=200] {1000*(5 - 2 * ( (6 * 0.01280755 * x)^(1/3) ))};
			
			\draw[markgrijs]       (axis cs:165, 334) -- (axis cs:165, 304) node[below, align=center, markgrijs] {\footnotesize 1};
			\draw[markgrijs]       (axis cs:168+1, 337) -- (axis cs:168+5, 340) node[right, align=center, markgrijs] {\footnotesize 2};
			\draw[markgrijs]       (axis cs:168+1, 327) -- (axis cs:168+7, 315) node[right, align=center, markgrijs] {\footnotesize 3};
			\draw[markgrijs]       (axis cs:168+1, 310) -- (axis cs:168+9, 290) node[right, align=center, markgrijs] {\footnotesize 4};
			
		\end{axis}
	\end{tikzpicture}
	\caption[Right-angle prism mirror geometry and dependence of mirror end facet length on polishing time]{Right-angle prism mirror geometry (left) and dependence of mirror end facet length on polishing time (right). The schematic shows the substrate dimensions in millimetres, as well as the end facet length $a$ and adjacent length $b$. The red volume $V$ represents the amount of removed material, which can be expressed in terms of the polishing depth $z$ or facet length $a$. On the right, each data point corresponds to a polished substrate (mirrors suitable for assembly are numbered 1-4). The solid line represents a fit based on Preston’s law, which assumes a constant material removal rate ($\dot{V}=\SI{0.0128}{\milli\meter\cubed\per\second}$).}
	\label{polishtime} 
\end{figure}

\begin{figure}[t]
	\centering
	\begin{tikzpicture}[remember picture]
		\node[anchor=south west,inner sep=0] (image1) at (0,0) {\includegraphics[width=0.97\textwidth]{{"mirror"}.pdf}};
		
		% left four
		\begin{scope}[x={(image1.south east)},y={(image1.north west)}]
			% draw a grid
			%\draw[help lines,xstep=.1,ystep=.1,overlay] (0,0) grid (1,1);
			% draw ticks
			%\foreach \x in {0,1,...,9} { \node [anchor=north,overlay] at (\x/10,0) {0.\x}; }
			%\foreach \y in {0,1,...,9} { \node [anchor=east,overlay] at (0,\y/10) {0.\y}; }
			
			% labels
			%\node[markgrijs, align=center, rotate=0] at (0.0,0.9) {\footnotesize\sc (a)};
			%\node[markgrijs, align=center, rotate=0] at (0.35,0.9) {\footnotesize\sc (b)};
			%\node[markgrijs, align=center, rotate=0] at (0.7,0.9) {\footnotesize\sc (c)};
			\node[markgrijs, align=center, rotate=0] at (0.10,0.555) {\footnotesize\sc (a)};
			\node[markgrijs, align=center, rotate=0] at (0.10+0.206,0.555) {\footnotesize\sc (b)};
			\node[markgrijs, align=center, rotate=0] at (0.10+0.412,0.555) {\footnotesize\sc (c)};
			
			\node[markgrijs, align=center, rotate=0] at (0.10,0.055) {\footnotesize\sc (d)};
			\node[markgrijs, align=center, rotate=0] at (0.10+0.206,0.055) {\footnotesize\sc (e)};
			\node[markgrijs, align=center, rotate=0] at (0.10+0.412,0.055) {\footnotesize\sc (f)};
			
			\node[markgrijs, align=center, rotate=0] at (0.815,0.055) {\footnotesize\sc (g)};
			
			
			\node[black, align=center, rotate=0] at (0.514,0.15) {\footnotesize\sc \SI{310}{\micro\meter}};
			\node[black, align=center, rotate=0] at (0.514,0.45) {\footnotesize\sc \SI{344}{\micro\meter}};
			\node[black, align=center, rotate=90] at (0.59,0.29) {\footnotesize\sc \SI{337}{\micro\meter}};
			\node[black, align=center, rotate=90] at (0.44,0.29) {\footnotesize\sc \SI{327}{\micro\meter}};
			
			\node[black, align=center, rotate=0] at (0.514-0.207-0.005,0.15+0.05) {\footnotesize\sc 4};
			\node[black, align=center, rotate=0] at (0.514-0.207-0.005,0.45-0.05) {\footnotesize\sc 1};
			\node[black, align=center, rotate=0] at (0.59-0.207-0.03,0.295) {\footnotesize\sc 2};
			\node[black, align=center, rotate=0] at (0.44-0.207+0.02,0.295) {\footnotesize\sc 3};
			
		\end{scope}
		
	\end{tikzpicture}
	\caption[Fabrication and assembly process of the pyramid MOT mirror system]{Fabrication and assembly process of the pyramid MOT mirror system. (\textsc{a}) Brass collet holding the aluminium mirror mount during machining. (\textsc{b}) Completed mirror holder, which is machined from a gold-coated, aluminium substrate mirror. The gold coating is protected with First Contact during fabrication of the mount. (\textsc{c}) A right-angle prism mirror after its adjacent corners at the base have been polished off. (\textsc{d}) Fully assembled pyramid MOT mirror system. (\textsc{e}) Microscope image showing a top view of the mirrors. The label numbers correspond to those in \autoref{polishtime}. (\textsc{f}) Magnified microscope view of the central aperture, with facet lengths indicated in micrometres. (\textsc{g}) Alignment tool used to position the mirrors within the mount during the glueing process. The alignment is performed via top-view microscope imaging.}
	\label{mirror}
\end{figure}

\subsection{MOT detection}

\begin{figure}[t]
	\centering
	\begin{tikzpicture}[remember picture]
		\node[anchor=south west,inner sep=0] (image1) at (0,0) {\includegraphics[width=1\textwidth]{{"pyramidtest"}.pdf}};
		
		% left four
		\begin{scope}[x={(image1.south east)},y={(image1.north west)}]
			% draw a grid
			%\draw[help lines,xstep=.1,ystep=.1,overlay] (0,0) grid (1,1);
			% draw ticks
			%\foreach \x in {0,1,...,9} { \node [anchor=north,overlay] at (\x/10,0) {0.\x}; }
			%\foreach \y in {0,1,...,9} { \node [anchor=east,overlay] at (0,\y/10) {0.\y}; }
			
			% labels
			\node[markgrijs, align=center, rotate=0] at (0.05,0.95) {\footnotesize\sc (a)};
			\node[markgrijs, align=center, rotate=0] at (0.57,0.95) {\footnotesize\sc (b)};
			\node[markgrijs, align=center, rotate=0] at (0.05,0.42) {\footnotesize\sc (c)};
			\node[markgrijs, align=center, rotate=0] at (0.57,0.42) {\footnotesize\sc (d)};
			
		\end{scope}
		
	\end{tikzpicture}
	\caption[Testbed setup for evaluating the standalone performance of a pyramid MOT before integration with a dipole trap and Fabry–Pérot resonator]{Testbed setup for evaluating the standalone performance of a pyramid MOT before integration with a dipole trap and Fabry–Pérot resonator. (\textsc{a}) A PTFE holder secures an aluminium mount containing four mirrors arranged in a pyramidal configuration. A cooling beam illuminates the entire pyramid, while a counter-propagating beam plugs the hole at the apex to balance the cooling forces and enable MOT formation. (\textsc{b}) Rear view of the mount, showing the gold-coated surface, which helps reflect additional cooling beams to create an optical molasses at the dipole trap position. (\textsc{c}) The stainless steel mounting finger supports the PTFE holder and aluminium mount, keeping the assembly in place while allowing optical access for the cooling beams. (\textsc{d}) Full assembly inside the glass cell, with anti-Helmholtz coils affixed to a 3D-printed frame, generating the quadrupole magnetic field required for MOT operation.}
	\label{pyramidtest} 
\end{figure}


\end{document}
